\PassOptionsToPackage{expansion=false,protrusion=false}{microtype}
\documentclass[sigconf,nonacm]{acmart}
\usepackage{booktabs}

\AtBeginDocument{%
  \providecommand\BibTeX{{\normalfont B\kern-0.5em{\scshape i\kern-0.25em b}\kern-0.25em T\kern-0.1667em\lower0.7ex\hbox{E}\kern-0.125emX}}}

\copyrightyear{2026}
\acmYear{2026}
\setcopyright{none}
\settopmatter{printacmref=false}
\renewcommand\footnotetextcopyrightpermission[1]{}
\pagestyle{plain}

\begin{document}

\title{PAnGEA: Grounding Free-Text Traffic Descriptions in Packet-Level Evidence}
\subtitle{Protocols \& Attacks in Networks: Grounding Evidence \& Assessment}

\author{Author Name}
\affiliation{%
  \institution{Institution Name}
  \city{City}
  \country{Country}
}
\email{author@example.edu}

\renewcommand{\shortauthors}{Author Name}

\begin{abstract}
Descriptions of network traffic -- whether produced by human analysts or by AI
systems -- can misrepresent the traffic they describe: a description can assert
activity the traffic does not support, or omit activity the traffic clearly
shows. We present PAnGEA, a system that takes a free-text traffic description
(an attack narrative, a protocol behavior description, or an ordinary traffic
summary) together with a packet capture (PCAP), decomposes the description into
atomic, individually checkable claims, and determines for each claim whether it
is \emph{grounded} in the traffic, \emph{unsupported} by it, or \emph{uncertain}
and worth further investigation. A complementary check flags significant traffic
activity with no corresponding claim in the description at all. Unlike existing
tools, which either extract structured features from a PCAP without reference to
any accompanying narrative (Zeek, NetworkMiner) or generate a narrative directly
from evidence (recent LLM-based incident-report generators), PAnGEA verifies an
\emph{already-existing} narrative against evidence, at claim-level granularity,
with an explicit accounting of unsupported and unaccounted-for content. We
describe PAnGEA's architecture, the design principles that emerged from
extensive iterative testing against real packet captures, a set of concrete
failure modes discovered and corrected during development, and an evaluation
methodology -- built around claim-level precision/recall against human-labeled
ground truth and established coreference-resolution metrics for actor
attribution -- for measuring the system's accuracy going forward.
\end{abstract}

\maketitle

% ============================================================
\section{Introduction}
% ============================================================

Network traffic captures are large, low-level, and difficult for a human analyst
to summarize completely and correctly. In the era of AI-generated incident
reports, drafts, and traffic summaries, a second, related risk has emerged:
a large language model (LLM) asked to describe a capture can hallucinate
plausible-sounding but unsupported activity, or omit activity that a purely
structural search would surface immediately. Both directions of error --
description asserting something the traffic does not show, and traffic showing
something the description does not mention -- can mislead an analyst who trusts
the description without independently re-examining the capture.

This paper presents PAnGEA (Protocols \& Attacks in Networks: Grounding
Evidence \& Assessment), a system whose task is not to \emph{generate} a
description of a capture, but to \emph{verify} a description that already
exists -- written by a human analyst, produced by another AI system, or
supplied as a protocol specification -- against the packet-level evidence in
a PCAP. PAnGEA decomposes the description into a sequence of atomic,
individually checkable claims and determines, for each claim, whether it is:

\begin{itemize}
  \item \textbf{Grounded} -- directly supported by structural evidence in the traffic;
  \item \textbf{Unsupported} -- a false positive, asserting activity the traffic does not show;
  \item \textbf{Uncertain} -- suggestive but not conclusively confirmable from the available evidence, and worth deeper investigation.
\end{itemize}

A complementary, symmetric check flags \emph{unaccounted evidence}: traffic
activity significant enough that a complete description should have mentioned
it, but did not (a false negative at the level of the narrative).

\subsection{Why not simply ask an LLM to check the whole capture?}

A natural alternative is to hand an LLM the raw PCAP and the description
together and ask it to judge consistency directly. Informal testing during this
project's development found that for small captures, this direct approach can
perform comparably to, or in some respects better than, the structured pipeline
described here -- an LLM reading raw packets directly retains information (exact
per-packet content, chronological ordering) that a compressed, aggregate
representation can lose. The structured pipeline exists to address a problem
this direct approach cannot: real packet captures routinely contain many
thousands to tens of thousands of packets, far exceeding what any current model
can accept as input at reasonable cost, and this scale problem does not
diminish as models' context windows grow, since capture sizes grow with them.
PAnGEA's Preprocessing stage (\S\ref{sec:preprocessing}) compresses a capture to
a budget-compliant representation without discarding the specific evidence a
claim-level check would need, making the same verification approach usable on
captures too large to hand to a model directly.

\subsection{Contributions}

\begin{itemize}
  \item A pipeline architecture that separates claim extraction from evidence access by construction, reducing the risk of a model's textual bias contaminating its own evidence-grounding judgment (\S\ref{sec:architecture}).
  \item A hybrid grounding strategy that resolves structural claims algorithmically where possible and defers to an LLM only where hand-written rules provably fail to generalize, informed by concrete failure cases discovered during development (\S\ref{sec:evidence-matching}).
  \item A protocol-agnostic structural anomaly detector for identity-fanout patterns (one low-level identifier associated with an anomalous number of distinct counterparts, relative to a capture's own baseline), generalized from an initial attack-specific version after that version was shown not to generalize (\S\ref{sec:anomaly-detection}).
  \item A joint actor-resolution mechanism that resolves all described actors in a single coordinated step rather than independently, addressing a concrete failure mode -- semantically incompatible roles resolving to the same network address -- observed repeatedly under an independent-resolution design (\S\ref{sec:joint-resolution}).
  \item An evaluation methodology, grounded in established techniques from summarization evaluation and coreference resolution, for measuring claim-level and actor-level accuracy against human-labeled ground truth (\S\ref{sec:evaluation-methodology}).
\end{itemize}

This paper describes the system's design, architecture, and the evaluation
methodology built to measure it. Quantitative results against a labeled dataset
are the subject of ongoing work and are not presented here.

% ============================================================
\section{Related Work}
% ============================================================

\textbf{Network traffic feature extraction.} Zeek~\cite{zeek} and
NetworkMiner~\cite{networkminer} both transform a raw PCAP into structured,
higher-level output -- Zeek into protocol-specific logs (\texttt{conn.log},
\texttt{http.log}, \texttt{dns.log}), NetworkMiner into host-centric profiles
(hostnames, operating systems, extracted credentials and files). Neither tool
accepts a free-text description as input, and neither produces a judgment about
whether a given textual claim is supported by the traffic. These tools address
a complementary problem -- extracting structured signal from a capture -- and
are candidate upstream sources of richer per-packet evidence for a system such
as PAnGEA, rather than substitutes for the verification task PAnGEA performs.

\textbf{Kill-chain and attack-narrative mapping.} Established practice in
network forensics maps sequences of packet events onto tactical frameworks such
as MITRE ATT\&CK~\cite{mitre_attack} (e.g., a SYN sweep mapped to
Reconnaissance, a periodic encoded beacon mapped to Command and Control).
PAnGEA's design was informed by this practice, particularly the observation
that treating a description as an unstructured flat list of steps discards
information a phase-aware structure would preserve.

\textbf{LLM-based incident narrative generation.} A recent line of tools uses
an LLM to generate a narrative report \emph{from} evidence: converting parsed
network logs into an attack storyline, or converting structured detection
findings into a written incident report. These systems operate in the direction
opposite to PAnGEA: evidence to narrative, rather than narrative to
verification against evidence. None of the systems examined during this
project's design phase perform claim-level decomposition of an
\emph{already-existing} narrative followed by independent verification of each
claim against packet evidence, with explicit false-positive/false-negative
accounting -- which is PAnGEA's specific contribution.

\textbf{Text-generation faithfulness evaluation.} PAnGEA's evaluation
methodology draws directly on two established lines of work outside the
network-security literature. The Pyramid method~\cite{nenkova2004pyramid}
evaluates summarization content selection using weighted gold content units,
where weight is derived from how many independent human summaries mention a
given unit -- the direct precedent for PAnGEA's severity-weighted gold claims
(\S\ref{sec:evaluation-methodology}). FActScore~\cite{min2023factscore}
evaluates long-form factual precision by decomposing generated text into atomic
facts and verifying each independently against a knowledge source -- the
precedent for PAnGEA's claim-level grounding methodology, applied here to
packet-level evidence rather than a text corpus.

\textbf{Coreference resolution evaluation.}
\label{sec:relwork:coref}
PAnGEA's actor-resolution accuracy
is measured using three established coreference-scoring metrics -- MUC
\cite{vilain1995muc}, B\textsuperscript{3}~\cite{bagga1998bcubed}, and
CEAF~\cite{luo2005ceaf} -- reframing "does this traffic description refer to
the same real-world host under two different names" as a coreference-clustering
problem. Each metric has a distinct, well-documented bias (MUC undercounts
singleton clusters; B\textsuperscript{3} can score an unmatched system mention
as trivially correct; CEAF's optimal-alignment step discards partial credit for
clusters that lose that alignment~\cite{cai2010twinless}), motivating PAnGEA's
choice to report all three rather than rely on any single one.

% ============================================================
\section{System Architecture}
\label{sec:architecture}
% ============================================================

PAnGEA is organized as a four-stage pipeline (Table~\ref{tab:stages}),
plus a fifth, offline Evaluator component used to measure the pipeline's
own accuracy rather than to process any individual sample.

\begin{table*}[t]
\caption{Pipeline stages, at a glance.}
\label{tab:stages}
\begin{tabular}{p{2.6cm}p{3.6cm}p{3.6cm}p{2.2cm}}
\toprule
\textbf{Stage} & \textbf{Input} & \textbf{Output} & \textbf{Method} \\
\midrule
Preprocessing & Raw PCAP + description (filter only) & Reduced, budget-compliant packet representation & Algorithmic \\
Step Extractor & Description text only & Actors + atomic claims & LLM \\
Evidence Matcher & Claims/actors + reduced evidence & Resolved actors + grounded claims & Hybrid \\
Validator & Grounder output + evidence & Extraction/linkage fidelity flags & LLM \\
\bottomrule
\end{tabular}
\end{table*}

\subsection{Description-blindness as a structural constraint}
\label{sec:r2}

A component that can see both a claim and the evidence for it at the same time
risks resolving ambiguity in the evidence toward whatever the claim already
led it to expect -- a form of confirmation bias transposed onto a model. PAnGEA
enforces a strict separation: the Preprocessing stage and the Evidence Matcher's
core lookup logic never receive the free-text description directly. Only the
Step Extractor sees the description, and only its already-extracted, structured
output (actors and atomic claims) flows to the later stages. This separation is
enforced at the level of function signatures and module boundaries, not by
convention alone, so that a future contributor cannot silently reintroduce the
dependency.

\subsection{Algorithmic-first, LLM-first: a decision made per-component, not globally}
\label{sec:hybrid}

PAnGEA does not adopt a single strategy (\emph{"prefer algorithms"} or
\emph{"prefer LLMs"}) uniformly across components; the appropriate strategy was
found to differ by task, and settled empirically rather than by prior
assumption.

Claim grounding (matching an atomic claim's structured indicators --
protocol, port, flags -- against a fingerprint index built from the reduced
traffic) is \textbf{algorithmic-first}: a direct dictionary lookup handles the
majority of claims at negligible cost, with escalation to an LLM reserved for
claims with multiple structurally competing matches or no usable structured
indicators at all.

Actor resolution (mapping a described role, such as "attacker" or "victim
host," to a specific IP address) is \textbf{LLM-first}. An earlier,
hand-written version scored candidate IPs using per-role formulas (e.g., an
"attacker" scores on destination-address diversity); this could not generalize
to any role the formulas' authors had not anticipated, and was replaced by an
approach that hands a role-agnostic structural summary of the traffic to an
LLM for interpretation against whatever role label a given scenario actually
uses (\S\ref{sec:evidence-matching}).

\subsection{Structural enforcement over trusting self-reported status}
\label{sec:enforcement}

Where a model's self-reported conclusion can be checked mechanically against
what it actually produced, PAnGEA checks it, rather than trusting the reported
label. Concretely: a claim's grounding status is downgraded automatically from
\emph{grounded} to \emph{uncertain} if no concrete supporting evidence was
actually linked, regardless of what status field an underlying model call
returned. This rule exists because the failure it prevents was observed
directly during development: a model reported a \emph{grounded} status backed
only by a plausible narrative explanation, with an empty evidence list attached.

% ============================================================
\section{Preprocessing}
\label{sec:preprocessing}
% ============================================================

The Preprocessing stage has two sub-components. A \emph{Filter Generator}
produces a Berkeley Packet Filter (BPF) expression from the description alone
(respecting the description-blindness boundary of \S\ref{sec:r2} with respect
to the \emph{Reducer} that follows it, though not with respect to the raw
description itself, whose only role here is to narrow the traffic under
consideration before any claim-level reasoning occurs). A \emph{Reducer} then
compresses the (optionally filtered) capture into a token-budget-compliant
representation, built around a protocol-agnostic \emph{structural
fingerprinting} scheme: for each packet, a signature is computed from
protocol-specific structural state alone (e.g., TCP flags and destination port;
ARP opcode and sender hardware address) -- deliberately excluding source/
destination address from the signature itself, so that structurally identical
traffic from many distinct hosts collapses into a single retained
representative plus an aggregate count, rather than one retained item per host.

Two properties of this design were only found to matter through direct
testing against real captures, rather than being anticipated in advance. First,
an under-specified token budget can silently truncate a capture after a tiny
fraction of its packets have been processed, with the resulting distortion
propagating through every later stage; the Reducer therefore reports an
explicit budget-exceeded status alongside the exact packet at which truncation
occurred, rather than truncating silently. Second, an encoding scheme for the
fingerprint signature that concatenates structural fields with a plain
separator character is unsafe when a field can itself contain that separator
(a MAC address contains colons); PAnGEA encodes fingerprint keys as
unambiguous structured strings rather than delimiter-joined text.

% ============================================================
\section{Step Extraction}
\label{sec:step-extraction}
% ============================================================

The Step Extractor is a necessarily LLM-based component: decomposing free text
into atomic, self-contained claims and identifying the actors referenced within
it require semantic understanding with no viable algorithmic substitute. Its
output is a flat list of atomic claims, each with an optional structured
\emph{expected-indicator} set (protocol, port, flags, direction, volume and
cardinality pattern) used for the algorithmic grounding path described in
\S\ref{sec:evidence-matching}, plus a free-text fallback description used when a
claim does not reduce cleanly to structured fields.

Iterative testing against real, human-annotated scenarios surfaced several
recurring extraction-quality failure modes, addressed through explicit
prompt-level constraints together with a self-check pass the model is asked to
perform on its own draft output before returning it:

\begin{itemize}
  \item \textbf{Entity fragmentation}: the same real host, referred to differently at different points in a narrative (e.g., introduced by function early on, referred to generically later), extracted as two separate actors.
  \item \textbf{Entity omission}: named infrastructure (a gateway, an external destination) present in the text but absent from the extracted actor list, despite later claims depending on it.
  \item \textbf{Category errors}: an abstraction (traffic direction itself) extracted as though it were a discrete network actor.
\end{itemize}

These prompt-level mitigations produced a measurable, but not complete,
reduction in the underlying error rate on repeated extraction of the same
scenario -- motivating the claim-level and actor-level accuracy metrics of
\S\ref{sec:evaluation-methodology}, which are designed specifically to quantify
this kind of residual, probabilistic extraction error rather than treat it as
either solved or unsolved.

% ============================================================
\section{Evidence Matching}
\label{sec:evidence-matching}
% ============================================================

The Evidence Matcher brings the extracted claims/actors and the reduced
traffic evidence together for the first time, performing two related but
distinct tasks: grounding each claim against the evidence, and resolving each
described actor to a specific network address.

\subsection{Claim grounding}

For each claim with usable structured indicators, a direct lookup against the
fingerprint index (\S\ref{sec:preprocessing}) either finds no candidates (the
claim is marked \emph{unsupported}), exactly one candidate (the claim is
marked \emph{grounded}, subject to the enforcement rule of
\S\ref{sec:enforcement}), or multiple competing candidates (escalated to an
LLM adjudication step). Claims with no usable structured indicators escalate
directly. One concrete matching defect, found only through testing against a
real ARP-spoofing capture, illustrates the value of testing every component
against real, not only synthetic, traffic: an early version of the lookup
logic rejected any candidate with flag-type criteria whose protocol was not
TCP, on the reasoning that "flags" are a TCP-specific concept. ARP has a
directly analogous single-value field (opcode: request versus reply) that the
Step Extractor reasonably maps into the same generic indicator field; the
effect of the defect was that claims describing the central mechanism of an
entire attack scenario were unconditionally marked unsupported, independent of
what evidence actually existed.

\subsection{Actor resolution and identity-fanout anomaly detection}
\label{sec:anomaly-detection}

Actor resolution hands an LLM a compact, per-address structural summary of the
traffic -- packet counts, destination/source cardinality, the dominant
fingerprint patterns associated with each address, and any \emph{structural
anomalies} detected. An anomaly is a traffic pattern whose fan-in or fan-out
cardinality is a statistical outlier relative to other traffic of the
\emph{same protocol in the same capture} -- a relative, capture-local baseline
rather than a fixed threshold, so the mechanism does not require re-tuning
across captures of very different size and density.

An initial version of this detector was written narrowly around one specific
ARP-spoofing signature (a single hardware address claiming ownership, via
unsolicited replies, of more than one IP address) and would not have
generalized to any other attack pattern. It was redesigned so that the
detection \emph{mechanism} is fully protocol-agnostic, while the necessarily
protocol-specific knowledge -- which sub-patterns of a given protocol even
constitute a claim worth checking (an ARP \emph{reply} claiming an identity;
not an ARP \emph{request}, which is routine) -- is isolated into a small,
explicit, extensible lookup table, added to incrementally as evidence from
real captures justifies each entry, rather than populated speculatively in
advance. This redesign was validated against three properties directly: the
original attack signature is still detected; an unrelated, synthetic anomaly
in a different protocol is also detected with no protocol-specific code
change; and ordinary high-volume traffic from a ordinary busy server, which an
earlier version of the mechanism had flagged as anomalous, is not.

\subsection{Joint resolution}
\label{sec:joint-resolution}

An initial actor-resolution design resolved each actor independently, in a
separate model call per actor. Repeated testing against a real capture
surfaced a persistent failure mode: two roles that must, by construction, refer
to different hosts (an attacker and a victim) resolved to an identical IP
address. Investigation traced this to a combination of causes -- reliance on
traffic-volume signals that do not actually discriminate between roles (a busy
attacker relaying traffic can resemble a busy legitimate host structurally),
and, more fundamentally, the complete absence of any mechanism for one
actor-resolution call to know what another call had concluded. An intermediate
mitigation -- informing each independent call of the \emph{labels} of the other
roles being resolved, without their answers -- measurably encouraged
self-directed caution but did not eliminate the underlying inconsistency.

The current design resolves all actors in a single joint call, given an
explicit instruction not to assign the same address to two roles without
specific justification. This is a structural rather than heuristic fix: a
model cannot fail to account for an earlier assignment it made within the same
continuous response, in the way it could across genuinely separate calls. A
secondary benefit is cost: the shared structural-summary context, previously
duplicated once per independent per-actor call, is now included once per
sample. This design is recent relative to the rest of the pipeline and has not
yet been subjected to the same depth of real-capture validation as the
components described above; known open questions include possible
attention degradation as the number of actors resolved in a single call grows,
and sensitivity to the order in which actors are presented within the prompt.

% ============================================================
\section{Validator}
\label{sec:validator}
% ============================================================

The Validator performs quality assurance on the Evidence Matcher's output
without independently re-deriving claim grounding or actor resolution from
scratch. It checks two distinct fidelity properties: \emph{extraction
fidelity} (does the extracted claim list faithfully and completely represent
the source description, judged against the text alone) and \emph{linkage
fidelity} (does each claim's grounding status actually match what the linked
evidence supports). The Validator is implemented but has, at the time of
writing, been deliberately withheld from exercising on real pipeline runs
while the Evidence Matcher's own design was still under active revision, to
avoid conflating instability in the component being checked with instability
in the checker.

% ============================================================
\section{Evaluation Methodology}
\label{sec:evaluation-methodology}
% ============================================================

PAnGEA's accuracy is intended to be measured, not judged by inspection, and
the Evaluator component exists to make that measurement repeatable across
pipeline changes. This section describes the methodology; quantitative results
against a labeled dataset are ongoing work.

\subsection{Gold construction}

A gold sample pairs one PCAP with one free-text description and a
manually-constructed reference: a set of gold actors (each with a
textual reference into the description and a ground-truth IP address,
verified by direct inspection of the capture) and a set of gold atomic
claims (each with a severity tag, described below, a gold grounding status,
and the specific fingerprint keys in the capture that support it, again
identified by direct manual inspection). Constructing this reference
independently of the system under evaluation is essential: using the
system's own output, or another automated tool's output, to help build the
reference risks the reference inheriting exactly the biases the evaluation
is meant to detect.

\subsection{Severity-weighted claim recall}

Not every missed claim is equally consequential: a claim describing a
distinct attack mechanism not captured by any other claim is a materially
worse omission than a claim redundant with one already captured. Following
the Pyramid method's precedent of weighting gold content units rather than
counting them uniformly~\cite{nenkova2004pyramid}, each gold claim is tagged
\emph{critical} or \emph{minor} at construction time, and claim recall is
computed as a severity-weighted quantity rather than a simple count.

\subsection{Claim and actor alignment}

Comparing system output to gold requires deciding which system claim (or
actor) corresponds to which gold claim (or entity) -- alignment must be based
on textual content, not on system-internal identifiers, since a system-chosen
label (an arbitrary snake-case string) carries no guaranteed relationship to
any gold label at all, and aligning by a resolved IP address instead would let
a resolution-stage error masquerade as an extraction-stage success. PAnGEA
aligns by similarity between each system item's textual reference and each
gold item's textual reference.

\subsection{Actor-resolution accuracy as coreference resolution}

Whether the system correctly recognized that two different textual mentions
refer to the same real-world host (or incorrectly split one host into two, or
merged two hosts into one) is structurally identical to the coreference
resolution evaluation problem in computational linguistics, and PAnGEA reports
all three of the field's standard metrics -- MUC~\cite{vilain1995muc},
B\textsuperscript{3}~\cite{bagga1998bcubed}, and CEAF~\cite{luo2005ceaf} --
rather than any single one, since each has a distinct, well-documented blind
spot (\S\ref{sec:relwork:coref}) and the combination is standard practice in
that literature for exactly this reason.

\subsection{Compliance testing}

A subset of correctness properties can be checked without any gold reference
at all, and can therefore be checked on every sample the pipeline processes,
not only the labeled subset: a claim's grounding status must never be
\emph{grounded} with no linked evidence at all (\S\ref{sec:enforcement}); every
resolved actor address must belong to the set of addresses actually observed
in that sample's traffic; and any two actors resolved to an identical address
are flagged for review. These checks were added directly in response to
concrete violations observed during development, rather than anticipated
speculatively.

\bibliographystyle{ACM-Reference-Format}
\bibliography{references}

\end{document}
